\section{Sequential contraction and decision lower bounds}

The approximation and private-moment bounds describe the precision attainable by the welfare procedure. The complementary lower bounds connect moment matching to the information retained by a sequential private experiment, and then connect that information to welfare estimation and honest interval length. The argument accommodates independent public randomness and message rules that depend on the preceding private messages, as specified in \cref{obj:def:sequential-class}.

\subsection{Measurable densities for conditional experiments}

Sequential experiments require density versions that are measurable jointly in the conditioning variables and the next message. For the auxiliary result below, \(H\) denotes a measurable parameter space and \(Z\) a standard-Borel output space. The finite kernels \(K\) and \(L\) assign output measures to each parameter value, with \(L\) serving as the reference kernel. A jointly measurable Radon--Nikodym density \(f\) represents \(K\) relative to \(L\) throughout this family of conditional experiments. The following result supplies that representation under pointwise absolute continuity.

\begin{theoremv}{lem:measurable-kernel-radon-nikodym}[Jointly measurable Radon–Nikodym density]
Let \(H\) and \(Z\) be measurable spaces, and let \(K,L:H\to Z\) be kernels satisfying:
\begin{itemize}
\item \textbf{(Standard Borel target.)} \(Z\) is a standard-Borel space.
\item \textbf{(Finiteness.)} Both \(K\) and \(L\) are finite kernels.
\item \textbf{(Pointwise domination.)} \(K(h,\cdot)\ll L(h,\cdot)\) for every \(h\in H\).
\end{itemize}
Then there exists a jointly measurable function \(f:H\times Z\to[0,\infty]\) such that, for every \(h\in H\) and every measurable set \(E\subseteq Z\),
\[
K(h,E)=\int_E f(h,z)\,L(h,dz).
\]
\end{theoremv}

The standard-Borel restriction applies to the output space, while the parameter space remains an arbitrary measurable space. This distinction allows the conditioning parameter to collect the public seed, the preceding messages, and the contrast vector. Joint measurability makes the density suitable for integration over both public randomness and parameter priors. The representation holds at every parameter value, matching the all-history formulation of privacy in \cref{obj:def:sequential-class}.

\subsection{Higher-order contraction under sequential privacy}

Fix a sequential protocol \(Q\) from \cref{obj:def:sequential-class}. For a public-seed realization \(r\), the reference experiment is the conditional transcript law under the zero-contrast causal law \(G_0\). The symmetric causal family \(G_\theta\) is defined in \cref{obj:def:symmetric-law}. Write \leanref{sym:p^Q}{\ensuremath{p^Q(\theta,r)}} for its conditional transcript density relative to that reference. The privacy contraction scale is \leanref{sym:\beta}{\ensuremath{\beta}}, equal to \((e^\varepsilon-1)/(2d)\). Conditioning on the seed preserves the protocol's dependence on the entire preceding message history.

To compare experiments generated by different contrast priors, let \leanref{sym:\alpha}{\ensuremath{\alpha}} be a positive amplitude at most \(1/2\). The coordinate priors \leanref{sym:\nu_0}{\ensuremath{\nu_0}} and \leanref{sym:\nu_1(same)}{\ensuremath{\nu_1}} are probability laws on \([-\alpha,\alpha]\), chosen independently of the public seed. Drawing the coordinates independently from either prior and then running the protocol gives the joint seed-transcript mixtures \leanref{sym:\mathsf M_0^Q}{\ensuremath{\mathsf M_0^Q}} and \leanref{sym:\mathsf M_1^Q}{\ensuremath{\mathsf M_1^Q}}. Moment matching means that the two coordinate priors have equal moments below a specified order. The next result bounds the resulting transcript proximity through the corresponding higher-order density derivative.

\begin{lemmav}{lem:adaptive-moment-contraction}[Adaptive higher-order transcript contraction]
Let \(n,d\) be integers and let \(\varepsilon\) be a privacy budget. Suppose that:
\begin{itemize}
\item \textbf{(Parameter regime.)} \(n\ge2\), \(d\ge2\), and \(0<\varepsilon\le1\).
\item \textbf{(Sampling and randomness.)} The sampling scheme satisfies \cref{obj:ass:iid-people,obj:ass:independent-randomness}.
\item \textbf{(Sequential privacy.)} \(Q\in\mathfrak Q_{\mathrm{SI}}(n,d,\varepsilon)\), as specified in \cref{obj:def:sequential-class}.
\item \textbf{(Measurable kernel densities.)} For every measurable parameter space \(H\), standard-Borel output space \(Z\), and finite kernels \(K,L\) from \(H\) to \(Z\) satisfying \(K(h,\cdot)\ll L(h,\cdot)\) for every \(h\), there exists a jointly measurable \(f:H\times Z\to[0,\infty]\) such that, for every \(h\) and measurable \(E\subseteq Z\),
\[
K(h,E)=\int_E f(h,z)\,L(h,dz).
\]
\end{itemize}
Set the privacy contraction scale to
\[
\beta=\frac{e^\varepsilon-1}{2d}.
\]
There exists a jointly measurable real-valued density version \(p^Q(\theta,r)(\mathbf z)\), defined for all \(\theta\in\mathbb R^d\), public seeds \(r\), and transcripts \(\mathbf z\), with the following properties. For every \(\theta\in[-1/2,1/2]^d\) and every \(r\), it is nonnegative at every transcript and gives the conditional transcript law under \(G_\theta\) relative to the zero-contrast reference:
\[
\Pr_Q(d\mathbf z\mid r,G_\theta)
=
p^Q(\theta,r)(\mathbf z)\,
\Pr_Q(d\mathbf z\mid r,G_0).
\]
For every \(\theta,r,\mathbf z,j\), its dependence on coordinate \(\theta_j\), with the other coordinates fixed, is a polynomial on \(\mathbb R\) of degree at most \(n\). For every \(\theta\in[-1/2,1/2]^d\), every \(r,j\), and every integer \(1\le k\le n\),
\[
\int
\left|\partial_{\theta_j}^{\,k}p^Q(\theta,r)(\mathbf z)\right|
\,\Pr_Q(d\mathbf z\mid r,G_0)
\le
k!\sqrt{\binom nk}\,\beta^k.
\]
For every \(\theta\in\mathbb R^d\), every \(r,\mathbf z,j\), and every integer \(k>n\),
\[
\partial_{\theta_j}^{\,k}p^Q(\theta,r)(\mathbf z)=0.
\]

For every amplitude \(0<\alpha\le1/2\), let \(\nu_0,\nu_1\) be probability laws supported on \([-\alpha,\alpha]\), fixed independently of the public seed. Define their seed-transcript mixtures by
\[
\mathsf M_\ell^Q(dr,d\mathbf z)
=
\lambda(dr)
\int_{[-\alpha,\alpha]^d}
\Pr_Q(d\mathbf z\mid r,G_\theta)\,
\nu_\ell^{\otimes d}(d\theta),
\qquad \ell\in\{0,1\}.
\]
For every integer \(k\ge1\) such that
\[
\int u^q\,\nu_0(du)=\int u^q\,\nu_1(du),
\qquad q=0,\ldots,k-1,
\]
the following conclusions hold:
\[
\operatorname{TV}(\mathsf M_0^Q,\mathsf M_1^Q)
\le
\min\left\{1,\,
d\alpha^k\sqrt{\binom nk}\,\beta^k\right\},
\qquad 1\le k\le n,
\]
and
\[
\mathsf M_0^Q=\mathsf M_1^Q,
\qquad k>n.
\]
Moreover, every Borel measurable seed-transcript decision \(f(r,\mathbf z)\in[0,1]\), interpreted as a randomized test's acceptance probability, satisfies
\[
\left|
\int f(r,\mathbf z)\,\mathsf M_0^Q(dr,d\mathbf z)
-
\int f(r,\mathbf z)\,\mathsf M_1^Q(dr,d\mathbf z)
\right|
\le
\operatorname{TV}(\mathsf M_0^Q,\mathsf M_1^Q).
\]
\end{lemmav}

The measurable-density condition is supplied by \cref{obj:lem:measurable-kernel-radon-nikodym}. The density has a probability interpretation on the admissible contrast cube and a coordinatewise polynomial extension to all real contrast vectors. Its degree bound reflects the finite participant count, while its integrated derivative bound quantifies the sensitivity of the private transcript at each order.

The contraction calculation retains the adaptive history within each conditional experiment. Both the density representation and the derivative bound hold for every public-seed realization, so the mixture conclusion accommodates any independent seed law \(\lambda\). The prior independence condition fixes the same comparison across those seed realizations. For matched moments below order \(k\), the displayed bound combines the amplitude factor \(\alpha^k\), the privacy factor \(\beta^k\), the participant factor \(\sqrt{\binom nk}\), and the dimension factor \(d\). Matching beyond the coordinatewise polynomial degree gives exact equality of the mixtures.

The final clause translates transcript proximity into a bound on every Borel randomized testing decision. Thus the result controls the information available to an analyst after observing both the public seed and all private messages. Together with the moment-matching measures in \cref{obj:lem:cai-low-moment-duality}, it supplies the experiment-comparison component of the welfare lower bounds.

\subsection{From welfare separation to decision lower bounds}

Transcript proximity becomes a decision lower bound when the two priors place their welfare targets near different centers. For two priors supported on the causal model \(\mathcal V_d\) from \cref{obj:def:causal-model}, let \(N_0\) and \(N_1\) denote their joint seed-transcript mixture laws. Write \(v_0<v_1\) for the welfare centers and \(\Delta=v_1-v_0\) for their separation. The quantity \(\eta_0\) bounds the probability, under each prior, that welfare lies farther than \(\Delta/8\) from its center; \(\omega\) bounds the total-variation distance between \(N_0\) and \(N_1\).

The following result separates these two requirements: concentration concerns the random target induced by a prior, whereas proximity concerns the distribution of the observable private experiment. Its decision laws also include the independent analyst randomizer \(U\) from \cref{obj:ass:independent-randomness}. This permits randomized estimators and intervals within the same comparison.

\begin{lemmav}{lem:separated-value-mixtures}[Separated welfare mixture bounds]
Fix a sample size \(n\), a dimension \(d\), a sampling scheme satisfying \cref{obj:ass:iid-people,obj:ass:independent-randomness}, and a law-independent one-message sequential protocol \(Q\) with probability message kernels and standard-Borel seed and message spaces. Write \(\mathcal V_d\) for the full-data probability laws satisfying the declared causal model of uniform strata, fair assignment, consistency, and interior arm means. The welfare target is
\[
V(P)=\frac{1}{d}\sum_{j=1}^{d}
\max\{\mu_{0j}(P),\mu_{1j}(P)\},
\]
where \(\mu_{aj}(P)\) is the conditional mean of potential outcome \(Y^a\) in stratum \(j\).

Consider two probability priors on full-data laws, called the first and second priors. Let \(v_0<v_1\) be their respective welfare centers, and define their separation by
\[
\Delta=v_1-v_0.
\]
Let \(\eta_0,\omega\in\mathbb R\). Suppose:
\begin{itemize}
\item \textbf{(Causal support.)} Each prior assigns probability one to \(\mathcal V_d\).
\item \textbf{(Decision mixtures.)} Mixing the induced joint law of \((\mathbf Z,R,U)\) over either prior yields a probability measure. Let \(N_0\) and \(N_1\) be the respective joint seed-transcript marginals, namely the mixture laws of \((R,\mathbf Z)\).
\item \textbf{(Target concentration.)} Under the first and second priors, respectively,
\[
\Pr\{|V(P)-v_0|>\Delta/8\}\le\eta_0,
\qquad
\Pr\{|V(P)-v_1|>\Delta/8\}\le\eta_0.
\]
\item \textbf{(Transcript proximity.)} The joint seed-transcript mixtures satisfy
\[
\operatorname{TV}(N_0,N_1)\le\omega.
\]
\end{itemize}

Then every Borel real-valued estimator \(T(\mathbf Z,R,U)\) satisfies
\[
\sup_{P\in\mathcal V_d}
\mathbb E_{P,Q}\bigl[(T(\mathbf Z,R,U)-V(P))^2\bigr]
\ge
\frac{9\Delta^2}{128}(1-\omega-2\eta_0)_+,
\]
with squared-error expectations interpreted in the extended nonnegative reals.

Moreover, every connected closed interval \(I(\mathbf Z,R,U)\) with ordered Borel endpoints in \([1/4,3/4]\) and global coverage
\[
\Pr_{P,Q}\{V(P)\in I(\mathbf Z,R,U)\}\ge0.90
\qquad\text{for every }P\in\mathcal V_d
\]
satisfies
\[
\sup_{P\in\mathcal V_d}
\mathbb E_{P,Q}\operatorname{length}I(\mathbf Z,R,U)
\ge
\frac{3\Delta}{4}(0.80-2\eta_0-\omega)_+.
\]
Here \(\operatorname{length}I\) is the upper endpoint minus the lower endpoint, and \(x_+=\max\{x,0\}\).
\end{lemmav}

The squared-error bound has the scale of the squared welfare separation. Its positive-part factor records the remaining testing separation after accounting for transcript proximity and the two target-concentration failures. Causal support ensures that averaging over either prior compares decisions within the same model used by the worst-law risk. Extended nonnegative expectations include every Borel real-valued estimator in this reduction.

For intervals, honesty is a finite-sample coverage requirement at every law in \(\mathcal V_d\). Consequently, the coverage requirement applies throughout the support of both priors. The factor \(0.80\) reflects the two \(0.90\) coverage requirements, with target-concentration and transcript-proximity penalties appearing explicitly. Connectedness gives the geometric content of the length bound: an interval containing target values within \(\Delta/8\) of both centers must span at least \(3\Delta/4\). Ordered Borel endpoints make coverage and length measurable, and their range agrees with the welfare range in \cref{obj:def:causal-model}.

Applied to independent-coordinate priors on the symmetric causal family, \cref{obj:lem:adaptive-moment-contraction} supplies the transcript-proximity bound. The bounded-mean concentration result in \cref{obj:lem:bounded-mean-concentration} supplies the corresponding control of average coordinate contributions to welfare. The separation supplied by approximation duality can therefore enter both decision criteria through \cref{obj:lem:separated-value-mixtures}: squared-error risk scales with \(\Delta^2\), and expected connected interval length scales with \(\Delta\).

\subsection{A dimension-free calibration}

A further comparison supplies a lower bound uniform in the number of strata. It uses the same measurable-density property and the sampling and randomness conditions already required for sequential contraction. The decision criteria are the minimax squared-error risk and globally honest connected interval length defined in \cref{obj:def:risk,obj:def:honest-length}. The next result states their dimension-free dependence on the finite privacy resources.

\begin{lemmav}{lem:dimension-free-private-two-point}[Dimension-free private two-point lower bounds]
Assume the following measurable Radon--Nikodym property: for every measurable parameter space \(H\), every standard-Borel output space \(Z\), and every pair of finite kernels \(K,L\) from \(H\) to \(Z\) satisfying \(K(h,\cdot)\ll L(h,\cdot)\) for every \(h\in H\), there exists a jointly measurable density \(f:H\times Z\to[0,\infty]\) such that, for every \(h\in H\) and every measurable \(E\subseteq Z\),
\[
K(h,E)=\int_E f(h,z)\,L(h,dz).
\]
Then there exists a universal constant \(c_1>0\) such that, for every sample size \(n\in\mathbb N\), dimension \(d\in\mathbb N\), privacy budget \(\varepsilon\in\mathbb R\), and sampling scheme satisfying
\begin{itemize}
\item \textbf{(Resources.)} \(n\ge2\), \(d\ge2\), and \(0<\varepsilon\le1\).
\item \textbf{(Participant sampling.)} \cref{obj:ass:iid-people}.
\item \textbf{(Randomization.)} \cref{obj:ass:independent-randomness}.
\end{itemize}
the minimax squared-error risk \(\mathfrak R_{\mathcal C}(n,d,\varepsilon)\) and globally \(0.90\)-honest connected interval length \(\mathfrak H_{\mathcal C}(n,d,\varepsilon)\), defined in \cref{obj:def:risk,obj:def:honest-length} for that sampling scheme, satisfy, for every \(\mathcal C\in\{\mathrm{NI},\mathrm{SI}\}\),
\[
\mathfrak R_{\mathcal C}(n,d,\varepsilon)
\ge c_1\min\left\{1,\frac{1}{n\varepsilon^2}\right\},
\qquad
\mathfrak H_{\mathcal C}(n,d,\varepsilon)
\ge c_1\min\left\{1,\frac{1}{\sqrt n\,\varepsilon}\right\}.
\]
\end{lemmav}

The two bounds depend on the participant count and privacy budget through \(n\varepsilon^2\), uniformly over every admissible dimension. When this quantity is at most one, they give constant lower bounds; above one, they give inverse-resource squared-error and inverse-square-root-resource length bounds. Both conclusions cover the noninteractive and sequential protocol classes, including the independent public and analyst randomness permitted by their definitions.

This comparison supplies a finite-resource baseline alongside the dimension-dependent moment-matching bounds. In particular, it applies at fixed small dimensions and throughout the stated sample-size range. Its universal constant allows the lower-bound calibration to retain the same dependence on privacy resources across dimensions, complementing the finite-sample calibrations in \cref{obj:lem:finite-private-polynomial-calibration} and the two-cell result in \cref{obj:thm:two-cell-calibration}.
